\section{A bounded subsolution theorem}

A bounded density relative to the Monge--Amp\`ere measure of one bounded potential also admits a bounded exponential solution. This permits right-hand sides singular with respect to smooth volume. Regularizing the reference gives a moving dominating measure; the mixed density inequality and fixed-reference trace estimate of \cref{lem:nonclosed-mixed-density,lem:bounded-reference-trace} bound the normalizations when the density is allowed to vanish.

\Needspace{15\baselineskip}
\begin{theorem}\label{thm:bounded-subsolution}
Let $\theta$ be semipositive and big. Suppose
\[
 \mu=f\,\mathrm{MA}_\theta(\psi),\qquad
 \psi\in\PSH(X,\theta)\cap L^\infty(X),\qquad
 0\leq f\in L^\infty(\mathrm{MA}_\theta(\psi)),
\]
and $\mu(X)>0$. For every $\lambda>0$ there is a unique bounded $u\in\PSH(X,\theta)$ satisfying
\[
 \mathrm{MA}_\theta(u)=\mathrm{e}^{\lambda u}\mu.
\]

\end{theorem}
\begin{proof}
Write $\nu=\mathrm{MA}_\theta(\psi)$, a finite non-pluripolar measure.

\emph{Smooth strictly positive density.} First suppose $f$ is smooth and strictly positive. Choose a Hermitian $\Omega\geq\theta$ and smooth $\Omega$-plurisubharmonic $\psi_j\downarrow\psi$, using the bounded regularization theorem. The rooftops
\[
 q_j=P_\theta(\psi_j)
\]
are uniformly bounded and decrease to $\psi$: they lie between $\psi$ and $\psi_j$. Contact support and ordered contact in background $\Omega$ give
\[
 \mathrm{MA}_\theta(q_j)
 \leq\mathbf1_{\{q_j=\psi_j\}}\Omega_{\psi_j}^n.
\]
Each right side has a bounded density with respect to smooth volume. Moreover $\mathrm{MA}_\theta(q_j)(X)>0$ by individual noncollapse. The bounded big-background theorem therefore supplies $u_j$ with
\[
 \mathrm{MA}_\theta(u_j)=\mathrm{e}^{\lambda u_j}f\,
                                     \mathrm{MA}_\theta(q_j).
\]
Since $f$ is bounded above and away from zero, $q_j-C$ and $q_j+C$ are lower and upper barriers for a fixed $C$. More explicitly choose $C$ so that
$\mathrm{e}^{\lambda(q_j-C)}f\leq1\leq \mathrm{e}^{\lambda(q_j+C)}f$ for all $j$. Exponential order then gives
$q_j-C\leq u_j\leq q_j+C$.

Extract an $L^1$ limit $u$, still uniformly bounded. The moving domination
\[
 \mathrm{MA}_\theta(u_j)\leq C'\Omega_{\psi_j}^n
\]
and \cref{cor:common-barrier-continuity}, with reference zero, give convergence in capacity. Bounded capacity continuity with the quasi-continuous weights $\mathrm{e}^{\lambda u_j}f$ gives
\[
 \mathrm{MA}_\theta(u)=\mathrm{e}^{\lambda u}f\nu.
\]

\emph{Bounded measurable density.} Choose positive smooth $f_j$ with a common upper bound $L>0$, converging to $f$ in $L^1(\nu)$. Such an approximation follows from regularity of the finite measure, continuous approximation on large compact sets, and smooth uniform approximation of the continuous functions; truncation to a fixed bounded interval and addition of a positive constant tending to zero preserve the bounds and convergence. The first step gives bounded $u_j$ with
\[
 \mathrm{MA}_\theta(u_j)=\mathrm{e}^{\lambda u_j}f_j\nu.
\]
A fixed $\psi-C$ is a lower barrier, since
$\mathrm{e}^{\lambda(\psi-C)}f_j\leq1$ for all $j$. Hence $u_j\geq\psi-C$. The upper bound requires a separate argument because $\inf_Xf_j$ can tend to zero.

The construction in the first step also gives the mixed inequality
\begin{equation}\label{eq:bounded-subsolution-mixed}
 \theta_{u_j}\wedge\theta_\psi^{n-1}
 \geq \mathrm{e}^{\lambda u_j/n}f_j^{1/n}\nu.
\end{equation}
To see this without assuming a mixed inequality for arbitrary singular measures, use that step's approximants $u_{j,k}$ and $q_k$. Their top measures have smooth-volume densities and
\[
 \mathrm{MA}_\theta(u_{j,k})
     =\mathrm{e}^{\lambda u_{j,k}}f_j\mathrm{MA}_\theta(q_k).
\]
Apply \cref{lem:nonclosed-mixed-density} to $(u_{j,k},q_k)$. For this fixed $j$ both sequences are uniformly bounded and converge in capacity to $(u_j,\psi)$. Bounded weighted continuity passes the mixed inequality to the limit and yields \eqref{eq:bounded-subsolution-mixed}.

Put $M_j=\sup_Xu_j$, $v_j=u_j-M_j$ and $c_j=\mathrm{e}^{\lambda M_j}$. If the $M_j$ were unbounded above, extract $M_j\to+\infty$ and then $v_j\to v$ in $L^1(X)$, with $v$ a genuine normalized potential. By \cref{lem:fixed-npp-integrals},
$\mathrm{e}^{\lambda v_j/n}\to \mathrm{e}^{\lambda v/n}$ in $L^1(\nu)$. The inequalities
\[
 |f_j^{1/n}-f^{1/n}|^n\leq|f_j-f|,
 \qquad0\leq f_j\leq L,
\]
then give
\[
 \int_X \mathrm{e}^{\lambda v_j/n}f_j^{1/n}\,\mathrm{d}\nu
 \longrightarrow\int_X\mathrm{e}^{\lambda v/n}f^{1/n}\,\mathrm{d}\nu>0.
\]
The last integral is positive because $\mu(X)>0$ and $\nu$ gives no mass to the pole set of $v$. Integrate \eqref{eq:bounded-subsolution-mixed} and apply \cref{lem:bounded-reference-trace}. For some $a>0$ along the extraction, this gives
\[
 a c_j^{1/n}\leq
 \int_X\theta_{v_j}\wedge\theta_\psi^{n-1}\leq C_\psi,
\]
contradicting $M_j\to+\infty$. Thus $u_j$ is uniformly bounded above as well as below.

Extract an $L^1$ limit. Now $\mathrm{MA}_\theta(u_j)\leq C''\nu$, so \cref{cor:common-barrier-continuity} gives capacity and weak Monge--Amp\`ere convergence. The right sides converge in total variation: use the fixed-measure exponential convergence and
\[
 \int_X|\mathrm{e}^{\lambda u_j}f_j-\mathrm{e}^{\lambda u}f|\,\mathrm{d}\nu
 \leq C\|f_j-f\|_{L^1(\nu)}
       +\|f\|_\infty\int_X|\mathrm{e}^{\lambda u_j}-\mathrm{e}^{\lambda u}|\,\mathrm{d}\nu.
\]
Consequently $\mathrm{MA}_\theta(u)=\mathrm{e}^{\lambda u}\mu$. Uniqueness is bounded exponential order.
\end{proof}
